
tôt de celle des carrés correspondants (7.10.0), compte tenu de la transitivité des images inverses (IV 2.12).

\textbf{Remarque 7.10.3.} Lorsque \(g\) est propre, la commutativité de (7.10.1) implique aussitôt celle du diagramme suivant, qui s'en déduit par passage aux gradués associés :
\begin{equation}\tag{7.10.1 bis}
\begin{tikzcd}[column sep=large,row sep=large]
\Gr_\bullet(X_t) \arrow[r,"g_{t*}"] \arrow[d,"\sigma_{X\bullet}"'] &
\Gr_\bullet(Y_t) \arrow[d,"\sigma_{Y\bullet}"]\\
\Gr_\bullet(X_s) \arrow[r,"g_{s*}"] &
\Gr_\bullet(Y_s).
\end{tikzcd}
\end{equation}
De même, si \(X\) et \(Y\) sont à fibres régulières et séparés, la commutativité de (7.10.2) implique celle de
\begin{equation}\tag{7.10.2 bis}
\begin{tikzcd}[column sep=large,row sep=large]
\Gr^\bullet(Y_t) \arrow[r,"g_t^*"] \arrow[d,"\sigma_Y^\bullet"'] &
\Gr^\bullet(X_t) \arrow[d,"\sigma_X^\bullet"]\\
\Gr^\bullet(Y_s) \arrow[r,"g_s^*"] &
\Gr^\bullet(X_s).
\end{tikzcd}
\end{equation}
Lorsque \(X\) et \(Y\) sont de plus quasi-projectifs et lisses sur \(S\), on sait d'ailleurs que \(g_t^*\) et \(g_s^*\) sont compatibles aux filtrations par la codimension des supports, et on trouve un diagramme commutatif comme (7.10.2 bis), mais les \(\Gr^\bullet\) étant remplacés par les \(\Gr^\bullet_{\mathrm{top}}\) (cf. Exp. 0, 2.1). On peut aussi noter que si \(X\) et \(Y\) sont quasi-projectifs sur \(S\) à fibres régulières et séparés, alors \(f_t\) et \(f_s\) sont automatiquement des morphismes d'intersection complète (VIII 1.1), et on obtient un diagramme commutatif analogue à (7.10.1 bis), où les \(\Gr_\bullet\) sont remplacés par les \(\Gr^\bullet(\ )_{\mathbb Q}\), compte tenu du résultat de compatibilité (VIII 3.2).

